%! Characteristics of Rational Functions — Unit 4, Lesson 4.0 — Guided Notes
\documentclass[10pt]{article}
\usepackage{algebra2-article}
\usepackage{algebra2-boxes}
\newcommand{\TallMath}[1]{$\displaystyle #1\rule[-1.4em]{0pt}{3.2em}$}
\newcommand{\lgrid}[4]{%
  \draw[step=1, linegray!40, very thin] (#1,#3) grid (#2,#4);
  \draw[->, semithick, charcoal] (#1,0)--(#2,0) node[below right, font=\scriptsize]{$x$};
  \draw[->, semithick, charcoal] (0,#3)--(0,#4) node[left, font=\scriptsize]{$y$};}

\begin{document}

\pageheader{Unit 4, Lesson 4.0}{Guided Notes}
\namedateperiod

\vspace{0.08in}

\begin{objectivebox}
\textbf{By the end of this lesson, I will be able to\dots}
\begin{itemize}\setlength{\itemsep}{2pt}
  \item find a rational function's \emph{domain restrictions}, and tell a \emph{vertical asymptote} from a \emph{hole};
  \item write the \emph{equations} of the vertical and horizontal asymptotes, from a graph or an equation, and use the horizontal asymptote to state \emph{end behavior};
  \item read a rational graph's \emph{domain}, \emph{range}, \emph{intercepts}/\emph{zeros}, and \emph{increasing/decreasing} and \emph{positive/negative} intervals, and contrast it with a polynomial graph.
\end{itemize}
\end{objectivebox}

\vspace{0.05in}

\begin{vocabbox}
\small
Fill in each term as we build it together.
\par\vspace{2pt}   % \par is required: \termblanklong's \noindent is a no-op mid-paragraph
\termblanklong{Rational function}
\termblanklong{Domain restriction (excluded value)}
\termblanklong{Vertical asymptote}
\termblanklong{Horizontal asymptote}
\termblanklong{Hole (removable discontinuity)}
\end{vocabbox}

\begin{hookbox}
Here is the graph of $f(x) = \dfrac{x+2}{x-1}$.
\begin{center}
\begin{tikzpicture}[scale=0.40]
  \lgrid{-4}{5}{-4}{5}
  \draw[navy, dashed, semithick] (1,-4)--(1,5);
  \draw[navy, dashed, semithick] (-4,1)--(5,1);
  \node[navy, font=\scriptsize, above] at (1,5) {$x=1$};
  \node[navy, font=\scriptsize, below] at (4.3,1) {$y=1$};
  \begin{scope}
    \clip (-4,-4) rectangle (5,5);
    \draw[very thick, forest, domain=-4:0.4, samples=90, smooth]
      plot (\x,{((\x)+2)/((\x)-1)});
    \draw[very thick, forest, domain=1.75:5, samples=90, smooth]
      plot (\x,{((\x)+2)/((\x)-1)});
  \end{scope}
  \fill[charcoal] (-2,0) circle (3pt) node[above left, font=\scriptsize]{$(-2,0)$};
  \fill[charcoal] (0,-2) circle (3pt) node[right, font=\scriptsize]{$(0,-2)$};
\end{tikzpicture}
\end{center}
\begin{itemize}\setlength{\itemsep}{1pt}
  \item Every graph we have read this year --- lines, parabolas, polynomials --- was \emph{one unbroken curve}. This one comes in \textbf{two pieces}. What input is missing, and why? \writeline
  \item At both far ends the curve hugs the line $y=1$ but never lands on it. What is $f(x)$ doing as $x$ runs off to the left and to the right? \writeline
\end{itemize}
A rational function is a \emph{fraction} of polynomials, and the \textbf{denominator runs the show}:
it decides which inputs are thrown out, where the graph breaks, and what the graph settles down to at
the far ends. Today we learn to read all of it.
\end{hookbox}

\begin{notesbox}{1. What Makes a Function Rational --- and Where It Is Undefined}
A \textbf{rational function} is a \emph{quotient of two polynomials}:
\[
  R(x) = \frac{P(x)}{Q(x)}, \qquad Q(x) \neq 0 .
\]
Division by zero is undefined, so \emph{every} input that makes the denominator zero is thrown out of
the domain. Such an input is called a \textbf{domain restriction} or an \textbf{excluded value}.

\vspace{3pt}
\textbf{The one procedure behind this whole unit:} to find the restrictions, set the
\blank{3.6cm} equal to \blank{1.0cm} and solve.
\begin{itemize}[leftmargin=1.5em, itemsep=3pt]
  \item For $f(x)=\dfrac{x+2}{x-1}$: solve $x-1=0$, so $x=\blank{0.8cm}$ is excluded.
        Domain: \blank{5.4cm}.
  \item For $g(x)=\dfrac{5}{x^2-9}$: solve $x^2-9=0$, so $x=\blank{2.0cm}$ are excluded.
        Domain: \blank{5.4cm}.
  \item \textbf{Careful:} the \emph{numerator} is allowed to be zero --- that just makes
        $R(x)=\blank{0.8cm}$, which is an $x$-intercept (a \textbf{zero} of the function), not a
        restriction.
\end{itemize}

\vspace{3pt}
Because inputs are missing, a rational graph can be \textbf{discontinuous} --- it can have breaks.
This is the first function family of the year whose graph is not one connected piece. A break shows up
in exactly \textbf{two} ways, and the rest of the lesson is about telling them apart:
\begin{center}
\small
a \textbf{vertical asymptote} (the graph flies off near that input) \quad or \quad
a \textbf{hole} (a single missing point).
\end{center}
\end{notesbox}

\begin{notesbox}{2. Vertical Asymptotes: the Walls}
A \textbf{vertical asymptote} is a vertical line the graph approaches but \emph{never} touches or
crosses. Near it the outputs grow without bound --- exactly what Warm-Up item~2 showed.

\vspace{3pt}
\textbf{How to find it.} If $x=a$ is a domain restriction and the factor $(x-a)$ does \emph{not}
cancel, then $x=a$ is a vertical asymptote.
\begin{itemize}[leftmargin=1.5em, itemsep=3pt]
  \item Write the answer as an \textbf{equation of a line}, not just a number:
        \blank{1.8cm}, not ``$1$.''
  \item For the anchor $f(x)=\dfrac{x+2}{x-1}$, the vertical asymptote is \blank{2.0cm}.
  \item A graph can \emph{never} cross a vertical asymptote, because that input is
        \blank{4.4cm} --- there is no point to plot.
  \item The two branches on either side of a vertical asymptote are why the domain must be written
        as a \textbf{union}: for $f$, domain $=(-\infty,\ \blank{0.7cm}) \cup (\blank{0.7cm},\ \infty)$.
\end{itemize}
\end{notesbox}

\begin{notesbox}{3. Horizontal Asymptotes \& End Behavior: the Level}
A \textbf{horizontal asymptote} is a horizontal line the graph approaches as $x\to-\infty$ and
$x\to+\infty$. It \emph{is} the end-behavior statement for a rational function: instead of arms that
run off to $\pm\infty$ like a polynomial's, a rational graph's arms usually \textbf{flatten out
toward a level}.

\vspace{2pt}
For the anchor: as $x\to+\infty,\ f(x)\to\blank{0.8cm}$; \; as $x\to-\infty,\ f(x)\to\blank{0.8cm}$;
\; so the horizontal asymptote is \blank{2.0cm}.

\vspace{4pt}
\textbf{Finding it from the equation: compare the degrees.} Let $n=\deg$ of the numerator and
$m=\deg$ of the denominator.
\begin{center}
\renewcommand{\arraystretch}{1.35}
\begin{tabularx}{\linewidth}{@{}l l X@{}}
\hline
\textbf{Case} & \textbf{Horizontal asymptote} & \textbf{Example} \\
\hline
$n<m$ (bottom-heavy) & $y=\blank{1.0cm}$ & $\dfrac{4}{x^2+1}$: \; HA is \blank{1.8cm} \\
$n=m$ (tie) & $y=\dfrac{\text{lead coef.\ of }P}{\text{lead coef.\ of }Q}$ & $\dfrac{x+2}{x-1}$: \; HA is \blank{1.8cm} \\
$n>m$ (top-heavy) & \blank{2.6cm} & $\dfrac{x^2}{x+1}$: \; the outputs grow \blank{2.4cm} \\
\hline
\end{tabularx}
\end{center}

\vspace{2pt}
\textbf{Two warnings.}
\begin{itemize}[leftmargin=1.5em, itemsep=3pt]
  \item A horizontal asymptote describes the \emph{far ends} only. Unlike a vertical asymptote, a
        graph \textbf{may} cross its horizontal asymptote in the middle. (Check: $y=\dfrac{2x}{x^2+1}$
        has HA $y=0$, and at $x=0$ the output is $\blank{0.8cm}$ --- so it \emph{sits on} the line
        there.)
  \item In the top-heavy case there is \emph{no} horizontal asymptote at all --- the graph's end
        behavior looks like a \blank{3.0cm} instead.
\end{itemize}
\end{notesbox}

\begin{notesbox}{4. Holes: When a Restriction Cancels}
Sometimes a restriction does \emph{not} produce a wall. Look at
\[
  h(x) = \frac{x^2-4}{x-2} = \frac{(x-2)(x+2)}{x-2} = x+2, \qquad x \neq 2 .
\]
The factor $(x-2)$ cancels, so the graph is just the \emph{line} $y=x+2$ --- with the single point at
$x=2$ punched out. That missing point is a \textbf{hole}, also called a \textbf{removable
discontinuity}.
\begin{center}
\begin{tikzpicture}[scale=0.46]
  \lgrid{-4}{4}{-3}{7}
  \begin{scope}
    \clip (-4,-3) rectangle (4,7);
    \draw[very thick, forest] (-4,-2)--(4,6);
  \end{scope}
  \draw[forest, very thick, fill=white] (2,4) circle (4.2pt);
  \node[charcoal, font=\scriptsize, right] at (-1.2,6.2) {the hole};
  \draw[charcoal, thin, ->] (-0.1,6.0) -- (1.75,4.25);
\end{tikzpicture}
\end{center}
\begin{itemize}[leftmargin=1.5em, itemsep=3pt]
  \item \textbf{The cancel test.} Factor the top and bottom. A restriction whose factor
        \textbf{cancels} gives a \blank{1.6cm}; a restriction whose factor \textbf{stays} in the
        denominator gives a \blank{4.6cm}.
  \item \textbf{Find the hole's coordinates} by substituting the excluded value into the
        \emph{simplified} expression: here $x=2$ gives $y=\blank{0.8cm}$, so the hole is at
        $(\blank{0.6cm},\ \blank{0.6cm})$.
  \item \textbf{Cancelling never restores the input.} Even though the simplified form $x+2$ looks
        perfectly happy at $x=2$, the original function is still \blank{3.0cm} there. Always read
        restrictions off the \textbf{original} denominator.
  \item \textbf{Both at once.} For $k(x)=\dfrac{x-3}{x^2-9}=\dfrac{x-3}{(x-3)(x+3)}=\dfrac{1}{x+3}$,
        $x\neq 3$: the restriction $x=3$ cancels, so it is a \blank{1.4cm}; the restriction
        $x=-3$ does not cancel, so it is a \blank{4.0cm}.
\end{itemize}
\end{notesbox}

\begin{notesbox}{5. The Full Feature List --- and How Rational Differs from Polynomial}
Now read \emph{every} characteristic off the anchor $f(x)=\dfrac{x+2}{x-1}$.
\begin{center}
\begin{tikzpicture}[scale=0.52]
  \lgrid{-4}{6}{-5}{6}
  \draw[navy, dashed, semithick] (1,-5)--(1,6);
  \draw[navy, dashed, semithick] (-4,1)--(6,1);
  \node[navy, font=\scriptsize, above] at (1,6) {$x=1$};
  \node[navy, font=\scriptsize, below] at (5.0,1) {$y=1$};
  \begin{scope}
    \clip (-4,-5) rectangle (6,6);
    \draw[very thick, forest, domain=-4:0.5, samples=100, smooth]
      plot (\x,{((\x)+2)/((\x)-1)});
    \draw[very thick, forest, domain=1.6:6, samples=100, smooth]
      plot (\x,{((\x)+2)/((\x)-1)});
  \end{scope}
  \fill[charcoal] (-2,0) circle (2.6pt) node[above left, font=\scriptsize]{$(-2,0)$};
  \fill[charcoal] (0,-2) circle (2.6pt) node[right, font=\scriptsize]{$(0,-2)$};
\end{tikzpicture}
\end{center}
\begin{enumerate}[leftmargin=1.7em, itemsep=3pt]
  \item \textbf{Vertical asymptote:} \blank{1.8cm} \quad \textbf{Horizontal asymptote:}
        \blank{1.8cm}
  \item \textbf{Domain:} \blank{5.0cm} \quad \textbf{Range:} \blank{5.0cm} \\
        (The range excludes the horizontal-asymptote value because the graph never \emph{reaches}
        it.)
  \item \textbf{$x$-intercept / zero:} \blank{1.8cm} --- found by setting the \blank{2.4cm} equal to
        zero. \quad \textbf{$y$-intercept:} \blank{1.8cm} --- found by evaluating $f(0)$.
  \item \textbf{Increasing / decreasing:} the curve falls on \emph{both} branches, so $f$ is
        decreasing on \blank{2.6cm} and on \blank{2.6cm}. \emph{Never} write a single interval
        across a vertical asymptote --- the function is undefined there.
  \item \textbf{Positive / negative:} $f(x)>0$ on \blank{4.4cm}; \quad $f(x)<0$ on \blank{2.4cm}.
  \item \textbf{End behavior:} as $x\to\pm\infty$, $f(x)\to\blank{0.8cm}$.
\end{enumerate}

\vspace{4pt}
\textbf{Compare and contrast} (this is what makes rational functions a new family):
\begin{center}
\renewcommand{\arraystretch}{1.3}
\begin{tabularx}{\linewidth}{@{}l X X@{}}
\hline
 & \textbf{Polynomial (Unit 3)} & \textbf{Rational (Unit 4)} \\
\hline
Graph & smooth and \emph{continuous} & may be \blank{2.6cm} (breaks) \\
Domain & \emph{all} real numbers & all reals \blank{3.0cm} \\
Asymptotes & \blank{1.6cm} & vertical and/or horizontal \\
End behavior & arms run to $\pm\infty$ & arms usually flatten toward \blank{2.4cm} \\
\hline
\end{tabularx}
\end{center}
\end{notesbox}

\begin{practicebox}
The graph below is $r(x) = \dfrac{4}{x^2-4} = \dfrac{4}{(x-2)(x+2)}$. Read each characteristic and
write it on the line.
\begin{center}
\begin{tikzpicture}[scale=0.5]
  \lgrid{-5}{5}{-5}{5}
  \draw[navy, dashed, semithick] (-2,-5)--(-2,5);
  \draw[navy, dashed, semithick] (2,-5)--(2,5);
  \node[navy, font=\scriptsize, above] at (-2,5) {$x=-2$};
  \node[navy, font=\scriptsize, above] at (2,5) {$x=2$};
  \begin{scope}
    \clip (-5,-5) rectangle (5,5);
    \draw[very thick, forest, domain=-5:-2.19, samples=80, smooth] plot (\x,{4/((\x)*(\x)-4)});
    \draw[very thick, forest, domain=-1.789:1.789, samples=80, smooth] plot (\x,{4/((\x)*(\x)-4)});
    \draw[very thick, forest, domain=2.19:5, samples=80, smooth] plot (\x,{4/((\x)*(\x)-4)});
  \end{scope}
  \fill[charcoal] (0,-1) circle (2.6pt) node[below right, font=\scriptsize]{$(0,-1)$};
\end{tikzpicture}
\end{center}
\begin{enumerate}[leftmargin=1.7em, itemsep=3pt]
  \item Vertical asymptotes: \blank{1.4cm} and \blank{1.4cm}. \quad Domain: \blank{5.4cm}
  \item Horizontal asymptote: \blank{1.6cm}. \quad Which degree case is this
        ($n<m$, $n=m$, $n>m$)? \blank{1.6cm}
  \item $y$-intercept: \blank{1.6cm}. \quad Does $r$ have an $x$-intercept? \blank{1.2cm}
        Why or why not? \blank{5.0cm}
  \item How many pieces does this graph come in? \blank{1.0cm} \quad Is $r(x)$ positive or negative
        between the two asymptotes? \blank{1.8cm}
\end{enumerate}
\end{practicebox}

\end{document}
